\subsection{Data Manipulation Language}
\begin{frame}[fragile]{SELECT -- The Foundation}
  The \texttt{SELECT} statement retrieves data from one or more tables.
  It is the most frequently used SQL statement.

  \begin{block}{Full Syntax (clauses in writing order)}
\begin{lstlisting}[style=sqlstyle]
SELECT column_name(s)
FROM table_name
WHERE condition
GROUP BY column_name(s)
HAVING condition
ORDER BY column_name(s)
LIMIT number
\end{lstlisting}
  \end{block}

  \begin{itemize}
    \item \texttt{SELECT} and \texttt{FROM} are the only \textbf{mandatory}
          clauses; all others are optional.
    \item \texttt{GROUP BY} is used together with aggregate functions;
          \texttt{HAVING} filters \emph{groups} (like \texttt{WHERE} but
          for groups, not individual rows).
  \end{itemize}

  \begin{alertblock}{Execution Order (not writing order!)}
    \texttt{FROM} $\to$ \texttt{WHERE} $\to$ \texttt{GROUP BY} $\to$
    \texttt{HAVING} $\to$ \texttt{SELECT} $\to$ \texttt{ORDER BY} $\to$
    \texttt{LIMIT}\\[2pt]
    Understanding this order is essential: it explains, for example, why
    column aliases defined in \texttt{SELECT} cannot be referenced in
    \texttt{WHERE}.
  \end{alertblock}
\end{frame}

% ----------------------------------------------------------
\begin{frame}[fragile]{SELECT -- Basic Usage}
  \begin{block}{Examples}
\begin{lstlisting}[style=sqlstyle]
-- Select all columns from a table
SELECT * FROM brand;

-- Select specific columns
SELECT name, address FROM brand;

-- Qualify column names (needed with multiple tables)
SELECT brand.name, brand.city FROM brand;
\end{lstlisting}
  \end{block}

  \begin{itemize}
    \item Using \texttt{*} returns \textbf{all columns} -- convenient for
          exploration but avoid in production code: column order can change
          when the table is altered, and transferring unneeded columns wastes
          network bandwidth and memory.
    \item Qualifying with \texttt{table.column} notation prevents
          \textbf{ambiguity} when joining multiple tables that share column
          names (e.g.\ both have an \texttt{id} column).
  \end{itemize}

  \begin{examples}{Tip: Explore first, then be specific}
    When learning a new schema, start with \texttt{SELECT * FROM table LIMIT
    10} to understand its structure, then narrow your query to only the
    columns you actually need.
  \end{examples}
\end{frame}

% ----------------------------------------------------------
\begin{frame}[fragile]{SELECT DISTINCT}
  \begin{block}{Returning Only Unique Values}
\begin{lstlisting}[style=sqlstyle]
-- Return only unique values
SELECT DISTINCT country FROM brand;

-- Multiple columns: the combination must be unique
SELECT DISTINCT country, legal_form FROM brand;
\end{lstlisting}
  \end{block}

  \begin{itemize}
    \item Without \texttt{DISTINCT}, SQL returns \textbf{all rows including
          duplicates}.  This differs from pure relational algebra, which
          always produces sets (no duplicates).
    \item \texttt{DISTINCT} applies to the \textbf{entire row}, not just one
          column: in the second example, the pair
          \texttt{(country, legal\_form)} must be unique.
  \end{itemize}

  \begin{alertblock}{Performance Warning}
    \texttt{DISTINCT} forces the database to sort or hash the entire result
    set to eliminate duplicates.  On large tables this can be very expensive.
    Avoid it if your query is already guaranteed to return unique rows, or if
    you simply want to count (use \texttt{COUNT(DISTINCT column)} instead).
  \end{alertblock}
\end{frame}

% ----------------------------------------------------------
\begin{frame}[fragile]{SELECT with LIMIT / TOP}
  \begin{block}{Restricting the Number of Returned Rows}
\begin{lstlisting}[style=sqlstyle]
-- MySQL / MariaDB / PostgreSQL:
SELECT * FROM product LIMIT 10;

-- Skip first 20 rows, then return the next 10
SELECT * FROM product LIMIT 10 OFFSET 20;

-- SQL Server syntax:
SELECT TOP 10 * FROM product;
\end{lstlisting}
  \end{block}

  \begin{itemize}
    \item \texttt{LIMIT} is essential when exploring large tables -- without
          it a query on a table with millions of rows could transfer
          gigabytes of data.
    \item \texttt{OFFSET} enables \textbf{pagination}: to show page~3 with
          10 rows per page, use \texttt{LIMIT 10 OFFSET 20}.
  \end{itemize}

  \begin{alertblock}{Always combine LIMIT with ORDER BY}
    Without an explicit \texttt{ORDER BY} the row order is
    \textbf{undefined} -- the database may return a different set of rows
    each time.  \texttt{LIMIT} without \texttt{ORDER BY} is only useful for
    a quick look at the data, never for reproducible results.
  \end{alertblock}
\end{frame}

% ----------------------------------------------------------
\begin{frame}[fragile]{Aliases (AS)}
  \begin{block}{Renaming Columns, Tables, and Expressions}
\begin{lstlisting}[style=sqlstyle]
-- Column alias
SELECT best_before_date AS bbd FROM stock;

-- Table alias (useful in JOINs)
SELECT p.name, p.price FROM product AS p;

-- Expression alias
SELECT AVG(price) AS avg_price FROM product;
\end{lstlisting}
  \end{block}

  \begin{itemize}
    \item Aliases are \textbf{temporary}: they exist only for the duration
          of the query and do not change the underlying table.
    \item Table aliases reduce typing and dramatically improve readability
          in complex queries with multiple joins.
    \item Column aliases can be used in \texttt{ORDER BY} (MySQL allows this
          as an extension) but \textbf{not} in \texttt{WHERE} or
          \texttt{HAVING} -- because those clauses are evaluated before
          \texttt{SELECT} assigns the alias.
  \end{itemize}

  \begin{examples}{Shorthand: AS is optional}
    Most DBMSs allow omitting the keyword \texttt{AS}:
    \texttt{SELECT price p FROM product} is equivalent to
    \texttt{SELECT price AS p FROM product}.
    Using \texttt{AS} explicitly is recommended for clarity.
  \end{examples}
\end{frame}

% ----------------------------------------------------------
\begin{frame}[fragile]{Aggregate Functions}
  Aggregate functions perform a calculation on a \textbf{set of rows} and
  return a \textbf{single value}.  They are the building block of
  statistical queries.

  \begin{block}{Common Aggregate Functions}
\begin{lstlisting}[style=sqlstyle]
SELECT COUNT(*) FROM product;         -- count all rows
SELECT COUNT(price) FROM product;     -- count non-NULL values
SELECT MIN(price) FROM product;
SELECT MAX(price) FROM product;
SELECT SUM(price) FROM product;
SELECT AVG(price) AS average_price FROM product;
\end{lstlisting}
  \end{block}

  \begin{itemize}
    \item \texttt{COUNT(*)} counts \textbf{all rows} (including those with
          \texttt{NULL}); \texttt{COUNT(column)} counts only rows where
          the column is \textbf{not} \texttt{NULL}.
    \item Aggregate functions are most powerful when combined with
          \texttt{GROUP BY} to compute per-group statistics.
  \end{itemize}

  \begin{alertblock}{Cannot Use Aggregates in WHERE}
    Aggregate functions are computed \emph{after} \texttt{WHERE} filtering.
    Use \texttt{HAVING} to filter on aggregate results.
  \end{alertblock}
\end{frame}
